\documentclass[11pt,a4paper]{article}
\usepackage[utf8]{inputenc}
\usepackage[T1]{fontenc}
\usepackage[spanish]{babel}
\usepackage{geometry}
\geometry{margin=2.5cm}
\usepackage{amsmath,amsfonts,amssymb}
\usepackage{booktabs,longtable,multirow,array}
\usepackage{graphicx}
\usepackage{hyperref}
\usepackage{siunitx}
\usepackage{caption,subcaption}
\usepackage{xcolor}
\usepackage{fancyhdr}
\usepackage{titlesec}
\usepackage{enumitem}
\usepackage{tikz}
\usepackage{colortbl}
\usepackage{float}

\definecolor{wasaffblue}{RGB}{0,82,147}
\definecolor{wasaffgray}{RGB}{100,100,100}
\definecolor{wasafflight}{RGB}{230,240,250}
\definecolor{accentgreen}{RGB}{34,139,34}

\pagestyle{fancy}
\fancyhf{}
\fancyhead[L]{\small\textcolor{wasaffblue}{WASAFF Consulting — Propuesta Comercial}}
\fancyhead[R]{\small\textcolor{wasaffblue}{CMPC Cordillera — Kaeser}}
\fancyfoot[C]{\thepage}
\renewcommand{\headrulewidth}{0.4pt}

\titleformat{\section}{\normalfont\Large\bfseries\color{wasaffblue}}{\thesection}{1em}{}
\titleformat{\subsection}{\normalfont\large\bfseries\color{wasaffblue!80}}{\thesubsection}{1em}{}
\titleformat{\subsubsection}{\normalfont\normalsize\bfseries\color{wasaffgray}}{\thesubsubsection}{1em}{}

\setlist[itemize]{leftmargin=1.5em,itemsep=0.3em}
\setlist[enumerate]{leftmargin=1.5em,itemsep=0.3em}

\begin{document}

% ============================================================
% PORTADA
% ============================================================
\begin{titlepage}
\begin{tikzpicture}[remember picture,overlay]
  \fill[wasaffblue] (current page.north west) rectangle ([yshift=-4cm]current page.north east);
  \node[anchor=north west,xshift=2cm,yshift=-1cm] at (current page.north west) {\color{white}\Large\textbf{WASAFF CONSULTING}};
  \node[anchor=north east,xshift=-2cm,yshift=-1cm] at (current page.north east) {\color{white}\large Propuesta Comercial};
\end{tikzpicture}

\vspace*{5cm}

\begin{center}
{\Huge\bfseries\textcolor{wasaffblue}{Sistema de Recuperación de Calor}}\\[0.4cm]
{\Huge\bfseries\textcolor{wasaffblue}{del Aceite Lubricante de Compresores}}\\[0.4cm]
{\LARGE\textcolor{wasaffgray}{CMPC Papeles Cordillera S.A. — Puente Alto}}\\[1.5cm]

\rule{0.6\textwidth}{1pt}\\[0.8cm]

{\Large\textbf{Propuesta de Ingeniería de Detalle}}\\[0.3cm]
{\large Nivel de Servicio: \textbf{Nivel 2 — Modelación Numérica Avanzada}}\\[0.3cm]
{\large Modalidad: \textbf{Ejecución + Entrega de Informe Técnico}}\\[1.5cm]

\begin{tabular}{rl}
\textbf{Cliente final:} & CMPC Papeles Cordillera S.A. \\[0.3em]
\textbf{Intermediario:} & Kaeser Compresores Chile SpA \\[0.3em]
\textbf{Consultor:} & WASAFF Consulting \\[0.3em]
\textbf{Contacto:} & Felipe Wasaff, Físico Computacional, MSc Candidate \\[0.3em]
\textbf{Fecha:} & Mayo 2026 \\[0.3em]
\textbf{Validez:} & 60 días corridos \\
\end{tabular}

\vfill

\begin{tikzpicture}
\draw[wasaffblue,thick,rounded corners] (0,0) rectangle (10,1.2);
\node at (5,0.6) {\large\textbf{\textcolor{wasaffblue}{Inversión: CLP \$\,480.000.000 + IVA (Nivel 2 Estándar)}}};
\end{tikzpicture}

\vspace{1cm}
{\small Documento confidencial. Uso exclusivo del destinatario.}
\end{center}
\end{titlepage}

% ============================================================
% ÍNDICE
% ============================================================
\tableofcontents
\newpage

% ============================================================
% 1. CARTA DE PRESENTACIÓN
% ============================================================
\section{Carta de Presentación}

Estimados señores de \textbf{Kaeser Compresores Chile SpA} y \textbf{CMPC Papeles Cordillera S.A.}:

Por medio de la presente, WASAFF Consulting tiene el agrado de presentar nuestra propuesta técnica y comercial para el desarrollo de la \textbf{Ingeniería de Detalle del Sistema de Recuperación de Calor del Aceite Lubricante de Compresores} en la planta de CMPC Papeles Cordillera, ubicada en Puente Alto, Región Metropolitana.

\subsection{Nuestra propuesta de valor}

A diferencia de una propuesta convencional basada en catálogos y factores de seguridad empíricos, esta propuesta integra \textbf{modelación matemática avanzada de nivel de maestría}:

\begin{itemize}
  \item \textbf{Método $\varepsilon$-NTU} con propiedades termofísicas dependientes de temperatura para el dimensionamiento preciso del intercambiador.
  \item \textbf{Dinámica transitoria} mediante integración Runge-Kutta de orden 4--5 (RK45), cuantificando el comportamiento del sistema ante arranques, paradas y fallas.
  \item \textbf{Análisis de sensibilidad paramétrica} con tornado diagrams, identificando los parámetros críticos que impactan el dimensionamiento y el riesgo técnico.
\end{itemize}

Este enfoque permite dimensionar con \textbf{menor incertidumbre}, reducir el sobre-dimensionamiento innecesario y entregar al cliente un modelo computacional validado que puede evolucionar hacia un \textbf{gemelo digital} para monitoreo predictivo.

\subsection{Experiencia relevante directa}

WASAFF Consulting cuenta con \textbf{experiencia directa en proyectos de recuperación de calor con CMPC Papeles Cordillera S.A.} En años anteriores participé en el diseño de un sistema de recuperación de calor para esta misma planta, lo que me proporciona un conocimiento único de sus procesos, infraestructura y requisitos específicos.

Adicionalmente, cuento con más de 6 años de experiencia en docencia universitaria en áreas de termofluidos, mecánica y análisis numérico, y actualmente curso el \textbf{Magíster en Simulación Computacional} en la Pontificia Universidad Católica de Valparaíso, profundizando en métodos numéricos y modelación matemática aplicada.

\vspace{1cm}

\begin{flushright}
\textit{Atentamente,}\\[0.5cm]
\textbf{Felipe Wasaff}\\
Físico Computacional, MSc Candidate\\
Director Técnico, WASAFF Consulting\\
Tel: +56 9 4612 5682\\
Email: felipe.wasaff@uchile.cl
\end{flushright}

\newpage

% ============================================================
% 2. RESUMEN TÉCNICO
% ============================================================
\section{Resumen Técnico del Sistema}

\subsection{Contexto de la planta}

CMPC Papeles Cordillera S.A. es la primera planta instalada por el grupo CMPC (año 1920), ubicada en la comuna de Puente Alto, a 20 km al sudeste de Santiago. Comercializa alrededor de 320.000 toneladas al año de papel para packaging, utilizando aproximadamente un \textbf{85\% de fibra reciclada} en su proceso.

La planta cuenta con una \textbf{central cogeneradora} compuesta por una turbina de gas Rolls-Royce Trent 60 (50 MW) y una caldera de recuperación de calor (HRSG) NEM con capacidad de 60 t/h de vapor a 16 bar / $213^\circ$C. Este contexto energético es relevante porque el sistema de recuperación de calor de los compresores debe evaluarse en función del \textbf{balance energético global} de la planta.

\subsection{Descripción del sistema propuesto}

El sistema consiste en una red hidráulica cerrada que extrae calor del aceite lubricante de los compresores de tornillo Kaeser y lo transfiere a agua de proceso mediante un intercambiador de calor de placas.

\begin{table}[H]
\centering
\caption{Componentes principales del sistema}
\begin{tabular}{@{}lll@{}}
\toprule
\textbf{Componente} & \textbf{Especificación} & \textbf{Nota} \\
\midrule
Intercambiador de calor & Placas soldadas, SS316L & Dimensionado por $\varepsilon$-NTU \\
Bomba de recirculación & Centrífuga con VFD & Curva del fabricante integrada \\
Acumulador térmico & Tanque aislado & Inercia térmica del sistema \\
Red de tuberías & Acero al carbono SCH40 & Darcy-Weisbach + Colebrook \\
Instrumentación & RTD 4-hilos, caudalímetros & Para calibración y control \\
\bottomrule
\end{tabular}
\end{table}

\subsection{Nota crítica sobre datos de entrada}

Los cálculos preliminares incluidos en esta propuesta están basados en \textbf{datos de catálogo de compresores Kaeser}. La verificación in-situ de los parámetros críticos (temperatura de aceite, caudales, presiones) es obligatoria para garantizar la precisión del dimensionamiento. Se incluye en el alcance una \textbf{visita técnica a planta} para medición y validación de estos datos.

\newpage

% ============================================================
% 3. ALCANCE DE SERVICIOS
% ============================================================
\section{Alcance de Servicios}

\subsection{Nivel de servicio: Nivel 2 — Modelación Numérica Avanzada}

\begin{table}[H]
\centering
\caption{Matriz de niveles de servicio WASAFF}
\begin{tabular}{@{}lccc@{}}
\toprule
\textbf{Capacidad} & \textbf{Nivel 1} & \textbf{Nivel 2} & \textbf{Nivel 3} \\
 & \textbf{(<$\$$3M)} & \textbf{($\$$3M--$\$$8M)} & \textbf{(>$\$$8M)} \\
\midrule
Cálculos Excel/Python básico & \checkmark & \checkmark & \checkmark \\
Propiedades $T$-dependientes & --- & \checkmark & \checkmark \\
ODEs (RK45, BDF) & --- & \checkmark & \checkmark \\
Análisis de sensibilidad & --- & \checkmark & \checkmark \\
Integración con cogeneración & --- & \checkmark & \checkmark \\
FEM 2D (FEniCS) & --- & \checkmark* & \checkmark \\
CFD 3D (OpenFOAM) & --- & --- & \checkmark \\
Gemelo digital & --- & --- & \checkmark \\
\bottomrule
\end{tabular}
\end{table}

*Validación FEniCS 2D incluida como entregable opcional.

\subsection{Entregables incluidos}

\subsubsection{Módulo 1: Análisis Estacionario ($\varepsilon$-NTU + Darcy-Weisbach)}
\begin{itemize}
  \item Framework computacional orientado a objetos en Python
  \item Siete escenarios de operación (mínimo a máximo)
  \item Dimensionamiento del intercambiador con propiedades termofísicas dependientes de $T$
  \item Cálculo de red hidráulica con iteración Colebrook-White
  \item Curva del sistema y punto de operación de la bomba
\end{itemize}

\subsubsection{Módulo 2: Análisis Transitorio (RK45)}
\begin{itemize}
  \item Modelo de capacitancia térmica lumped
  \item Simulación de tres eventos: arranque en frío, pérdida de compresor, arranque de backup
  \item Cuantificación del tiempo de respuesta $\tau_{95\%}$
\end{itemize}

\subsubsection{Módulo 3: Análisis de Sensibilidad Paramétrica}
\begin{itemize}
  \item Barrido paramétrico $\pm 10\%$ sobre variables críticas
  \item Tornado diagrams de impacto
  \item Curvas de sensibilidad métrica vs. parámetro
\end{itemize}

\subsubsection{Módulo 4: Integración con Sistema de Cogeneración}
\begin{itemize}
  \item Análisis de balance energético con turbina Trent 60 + HRSG existente
  \item Identificación de sinergias y conflictos energéticos
  \item Evaluación de usos prioritarios del calor recuperado
\end{itemize}

\subsubsection{Módulo 5: Documentación}
\begin{itemize}
  \item Informe técnico completo en LaTeX/PDF ($\sim$20--25 páginas)
  \item Informe ejecutivo en LaTeX/PDF (3--5 páginas)
  \item Memoria de cálculo firmada
  \item Especificaciones técnicas para licitación de equipos
\end{itemize}

\newpage

% ============================================================
% 4. ESPECIFICACIONES TÉCNICAS PRELIMINARES
% ============================================================
\section{Especificaciones Técnicas Preliminares}

\subsection{Intercambiador de calor de placas}
\begin{table}[H]
\centering
\begin{tabular}{@{}ll@{}}
\toprule
\textbf{Parámetro} & \textbf{Especificación} \\
\midrule
Tipo & Placas soldadas, contracorriente \\
Material & Acero inoxidable AISI 316L \\
Área de transferencia & A definir por modelación (estimado 25--35 m$^2$) \\
Coeficiente global $U$ & $\geq$ 3,000 W/m$^2$K (diseño) \\
Efectividad de diseño & $\varepsilon \geq$ 0.85 \\
Presión máxima de trabajo & 10 bar \\
Temperatura máxima & 150$^\circ$C \\
Conexiones & Bridas ANSI 150\# \\
Sellado & Soldadura láser (sin juntas) \\
\bottomrule
\end{tabular}
\end{table}

\subsection{Bomba de recirculación}
\begin{table}[H]
\centering
\begin{tabular}{@{}ll@{}}
\toprule
\textbf{Parámetro} & \textbf{Especificación} \\
\midrule
Tipo & Centrífuga horizontal, acople directo \\
Material & Cuerpo hierro fundido, impulsor bronce \\
Caudal nominal & 8--12 m$^3$/h (rango operativo) \\
Cabeza nominal & 5 m.c.a. (incluye margen) \\
Potencia motor & 1.5 kW \\
Control & Variador de frecuencia (VFD) \\
\bottomrule
\end{tabular}
\end{table}

\subsection{Instrumentación propuesta}
\begin{table}[H]
\centering
\begin{tabular}{@{}lll@{}}
\toprule
\textbf{Tag} & \textbf{Instrumento} & \textbf{Función} \\
\midrule
TT-101 & RTD Pt100 4-hilos & Temp. aceite entrada IC \\
TT-102 & RTD Pt100 4-hilos & Temp. aceite salida IC \\
TT-103 & RTD Pt100 4-hilos & Temp. agua entrada IC \\
TT-104 & RTD Pt100 4-hilos & Temp. agua salida IC \\
FT-101 & Caudalímetro electromagnético & Caudal agua de proceso \\
PI-101 & Transmisor de presión & Presión bomba descarga \\
\bottomrule
\end{tabular}
\end{table}

\newpage

% ============================================================
% 5. INVERSIÓN Y CONDICIONES COMERCIALES
% ============================================================
\section{Inversión y Condiciones Comerciales}

\subsection{Estructura de precios}

\begin{table}[H]
\centering
\caption{Desglose de inversión Nivel 2 Estándar}
\begin{tabular}{@{}lrr@{}}
\toprule
\textbf{Concepto} & \textbf{HH} & \textbf{Valor CLP} \\
\midrule
\rowcolor{wasafflight}
\multicolumn{3}{l}{\textbf{1. Ingeniería y Modelación}} \\
Análisis estacionario ($\varepsilon$-NTU + hidráulica) & 48 & \$\,76.800.000 \\
Análisis transitorio (RK45) & 32 & \$\,51.200.000 \\
Análisis de sensibilidad paramétrica & 20 & \$\,32.000.000 \\
Validación numérica cruzada & 12 & \$\,19.200.000 \\
Integración con cogeneración & 16 & \$\,25.600.000 \\
\midrule
\rowcolor{wasafflight}
\multicolumn{3}{l}{\textbf{2. Documentación}} \\
Informe técnico completo (LaTeX/PDF) & 32 & \$\,51.200.000 \\
Informe ejecutivo (LaTeX/PDF) & 12 & \$\,19.200.000 \\
Especificaciones técnicas para licitación & 24 & \$\,38.400.000 \\
Memoria de cálculo firmada & 16 & \$\,25.600.000 \\
\midrule
\rowcolor{wasafflight}
\multicolumn{3}{l}{\textbf{3. Gestión y Coordinación}} \\
Reuniones técnicas (5 reuniones de 2h) & 20 & \$\,32.000.000 \\
Gestión de proyecto + Kaeser & 24 & \$\,38.400.000 \\
Visita técnica a planta (Puente Alto) & 16 & \$\,25.600.000 \\
Presentación de resultados al cliente & 12 & \$\,19.200.000 \\
\midrule
\rowcolor{wasafflight}
\multicolumn{3}{l}{\textbf{4. Post-venta}} \\
Soporte 30 días + capacitación & 16 & \$\,25.600.000 \\
\midrule
\textbf{Total Nivel 2 Estándar} & \textbf{300} & \textbf{\$\,480.000.000} \\
\bottomrule
\end{tabular}
\end{table}

\textbf{Tarifa horaria:} CLP \$\,160.000/hora (4,05 UF/hora a mayo 2026).

\subsection{Condiciones de pago}

\begin{table}[H]
\centering
\caption{Condiciones de pago}
\begin{tabular}{@{}lcl@{}}
\toprule
\textbf{Hitos} & \textbf{\%} & \textbf{Monto CLP} \\
\midrule
Anticipo contra firma + info base entregada & 25\% & \$\,120.000.000 \\
Entrega informe conceptual aprobado & 25\% & \$\,120.000.000 \\
Entrega Módulos 1--2 (estacionario + transitorio) & 25\% & \$\,120.000.000 \\
Entrega final + presentación & 20\% & \$\,96.000.000 \\
Cierre (30 días post-entrega, sin observaciones) & 5\% & \$\,24.000.000 \\
\midrule
\textbf{Total} & \textbf{100\%} & \textbf{\$\,480.000.000} \\
\bottomrule
\end{tabular}
\end{table}

\subsection{Notas comerciales}

\begin{enumerate}
  \item Precios en pesos chilenos (CLP), valores \textbf{sin IVA}.
  \item Tarifa indexada a UF. Al 18 de mayo de 2026: UF $\approx$ \$\,39.500.
  \item El alcance \textbf{no incluye} equipos físicos. Solo ingeniería de detalle y modelación.
  \item Traslados asumen ubicación en \textbf{Puente Alto, RM}. Otras zonas se cotizan aparte.
  \item Entregables digitales en repositorio Git + archivos PDF/CSV/PNG.
  \item Tiempo de respuesta a consultas: 48 horas hábiles.
\end{enumerate}

\newpage

% ============================================================
% 6. PLAZO Y METODOLOGÍA
% ============================================================
\section{Plazo de Ejecución y Metodología}

\subsection{Cronograma}

\begin{table}[H]
\centering
\caption{Cronograma de 6 semanas}
\begin{tabular}{@{}lcc@{}}
\toprule
\textbf{Actividad} & \textbf{Semanas} & \textbf{HH} \\
\midrule
Kick-off y revisión de datos & 1 & 16 \\
Módulo 1: Análisis estacionario & 2--3 & 64 \\
Módulo 2: Análisis transitorio & 3--4 & 32 \\
Módulo 3: Sensibilidad paramétrica & 4--5 & 32 \\
Análisis de integración con cogeneración & 4--5 & 32 \\
Redacción de informes & 5 & 64 \\
Revisión interna y validación & 5--6 & 28 \\
Presentación al cliente y cierre & 6 & 32 \\
\midrule
\textbf{Total} & \textbf{6 semanas} & \textbf{300 HH} \\
\bottomrule
\end{tabular}
\end{table}

\subsection{Stack tecnológico}

\begin{table}[H]
\centering
\begin{tabular}{@{}lll@{}}
\toprule
\textbf{Capa} & \textbf{Herramienta} & \textbf{Uso} \\
\midrule
Lenguaje & Python 3.12 & Computación científica \\
Arrays & NumPy 2.4 & Álgebra lineal \\
ODEs & SciPy 1.17 (RK45) & Transitorios \\
Auto-diff & JAX 0.10 & Sensibilidad \\
FEM 2D & FEniCS legacy & Validación \\
Visualización & Matplotlib 3.10 & Figuras \\
Datos & Pandas 3.0 & Tablas CSV \\
Documentación & LaTeX & Informes formales \\
Control de versiones & Git & Trazabilidad \\
\bottomrule
\end{tabular}
\end{table}

\newpage

% ============================================================
% 7. CONDICIONES GENERALES
% ============================================================
\section{Condiciones Generales}

\subsection{Propiedad intelectual}
\begin{itemize}
  \item El código fuente (Python, LaTeX) es propiedad de WASAFF Consulting y se licencia al cliente para uso exclusivo en este proyecto.
  \item El cliente puede usar, modificar y extender el código para uso interno en CMPC.
  \item La publicación de resultados requiere autorización previa de WASAFF Consulting.
\end{itemize}

\subsection{Confidencialidad}
\begin{itemize}
  \item Ambas partes mantienen confidencialidad de la información técnica y comercial intercambiada.
  \item WASAFF no divulgará datos de CMPC a terceros sin autorización escrita.
\end{itemize}

\subsection{Limitación de responsabilidad}
\begin{itemize}
  \item La responsabilidad se limita al alcance definido en esta propuesta.
  \item No se incluye responsabilidad por fallas de equipos fabricados por terceros.
  \item Garantía de entregables: 30 días desde la aceptación formal.
\end{itemize}

\subsection{Aceptación y cierre}
\begin{itemize}
  \item El cliente dispone de 10 días hábiles para revisar y aceptar cada entregable.
  \item La falta de observaciones dentro del plazo se considera aceptación tácita.
  \item El cierre se formaliza con un acta de entrega-aceptación firmada.
\end{itemize}

\vspace{2cm}

\begin{center}
\rule{0.6\textwidth}{0.5pt}\\[0.5cm]
\textbf{Agradecemos la oportunidad de presentar esta propuesta.}\\[0.3cm]
\textit{Quedamos atentos a sus comentarios y a coordinar una reunión}\\
\textit{para discutir los detalles técnicos y comerciales.}\\[1cm]
\textbf{WASAFF Consulting}\\
\texttt{www.wasaff.cl} — \texttt{felipe.wasaff@uchile.cl}\\
+56 9 4612 5682
\end{center}

\end{document}
